\clearpage
\SourceLeaf{390}
\fancyhead[L]{\small VERSCHIEDENE AUFSÄTZE.}
\pdfbookmark[0]{Verschiedene Aufsätze.}{verschiedene-aufsaetze}
\pdfbookmark[1]{Bestimmung der größten Ellipse. 1810.}{groesste-ellipse-1810}
\begin{center}
{\Large VERSCHIEDENE AUFSÄTZE.}\par\bigskip
\rule{.85\textwidth}{.3pt}\par
v.\ \textsc{Zach}. Monatliche Correspondenz für Erd- und Himmelskunde. August 1810.\par
\rule{.85\textwidth}{.3pt}\par\bigskip
BESTIMMUNG DER GRÖSSTEN ELLIPSE\par\medskip
DIE DIE VIER SEITEN EINES GEGEBENEN VIERECKS BERÜHRT.\par\bigskip
\rule{.15\textwidth}{.3pt}
\end{center}
\hspace*{\parindent}Die Lage aller Punkte in der Ebene des Vierecks bestimme ich durch Abszissen und Ordinaten. Die Abszissenachse und ihren Ursprung wähle ich zunächst ganz nach Belieben. Zur Bestimmung des Vierecks verwende ich anstelle seiner Eckpunkte die Schnittpunkte seiner Seiten mit den Senkrechten, die vom Ursprung der Abszissenachse auf diese Seiten gefällt werden. Diese Senkrechten seien \(a,a',a'',a'''\), und ihre Neigungen gegen die Abszissenachse seien \(A,A',A'',A'''\), sodass die Koordinaten der genannten vier Schnittpunkte lauten
\[
\begin{array}{ll}
a\cos A,&a\sin A\\
a'\cos A',&a'\sin A'\\
a''\cos A'',&a''\sin A''\\
a'''\cos A''',&a'''\sin A'''
\end{array}
\]
